% =====================================================================
% Section 4: Method — EvidenceFlow
% =====================================================================
%
% Status: Draft §4 Method v8 (2026-05-17)
%   v8: §4 opening rewritten as a Figure 1 walkthrough per advisor
%       feedback ("opening should be a brief description of the
%       method figure"). Body §4.1/§4.2/§4.3 untouched in this revision.
%   v7.5 (pre-opening): §4.3 reframed as Evidence-Aware Retrieval
%       Optimization, 4 equations, problem-then-solution framing.
%   v7.4: bodies stable, pre-opening checkpoint archived.
%
% Method name: \methodname (defined via \newcommand in main file)
% Forward reference: fig:overview (figure to be added in figures/)
%
% =====================================================================

\section{Method: \methodname{}}
\label{sec:method}

\begin{figure}[t]
\centering
\fbox{\parbox{0.92\linewidth}{\centering\vspace{1.5em}
\textit{Overview of \methodname{}.}\\[0.3em]
Indexing: corpus $\to$ fact-source substrate.\\
Retrieval: $G_q \to H_q \to \pi_q \to s_{\mathrm{ppr}}$.\\
Optimization: relation-grounded expansion + coverage-aware reranking $\to R_q$.
\vspace{1.5em}}}
\caption{Overview of \methodname{}. Indexing builds a fact-source
substrate that links each OpenIE fact to its source passage through
endpoint and source edges. Given a query, a query-local subgraph
$G_q$ is induced over this substrate, and a personalized PageRank on
the derived graph $H_q$ produces an evidence flow $\pi_q$ and a base
passage score $s_{\mathrm{ppr}}$. Two evidence-aware steps then
extend the candidate pool through relation-grounded expansion and
form the final ranked evidence list $R_q$ via coverage-aware
reranking.}
\label{fig:overview}
\end{figure}

Figure~\ref{fig:overview} gives an overview of \methodname{}. At
indexing time, the corpus is parsed into OpenIE facts whose head and
tail phrases are linked through two complementary edge types,
endpoint edges across passages and source edges within a passage,
yielding a fact-source substrate that is built once for $\mathcal{D}$
(\S\ref{sec:method:substrate}). At query time, the question first
selects a query-local subgraph $G_q$ over this substrate by combining
dense entry passages with the provenance of facts whose endpoints
match question-side anchors, after which a personalized PageRank on a
derived graph $H_q$ propagates mass through fact nodes and reads out
a passage score $s_{\mathrm{ppr}}$
(\S\ref{sec:method:retrieval}). The local ranking is then refined by
two evidence-aware steps over the same substrate: relation-grounded
expansion admits passages whose relation phrases match $q$ even when
they fall outside $G_q$, and coverage-aware reranking selects the
final list $R_q$ greedily under a noisy-OR coverage model over
question-side requirements (\S\ref{sec:method:enhancement}).

\subsection{Offline Indexing: A Fact-Source Substrate}
\label{sec:method:substrate}

Our indexing differs from the standard passage--phrase formulation used
by HippoRAG-style systems in a single deliberate move: instead of
treating the OpenIE graph as a structure over passages and entity
phrases, we promote each OpenIE fact to a graph node and tie every
fact to its source passage. The resulting object is a
\emph{fact-source} substrate, indexed once for the corpus, on which
the query-time propagation of \S\ref{sec:method:retrieval} runs.

\paragraph{Fact units with provenance.}
For each document $d \in \mathcal{D}$, a language model extracts
OpenIE facts $f=(h,r,t)$ whose head and tail $h, t$ are argument
phrases and $r$ is the relation phrase
\citep{angeli2015leveraging,gutierrez2025hipporag}. We retain every
fact as a first-class indexing unit and bind it to its source through
a provenance map $\mathrm{src}(f)=d$, yielding a global fact set
$\mathcal{F}$ with the per-document slice
$\mathcal{F}_d = \{f\in\mathcal{F} : \mathrm{src}(f)=d\}$. Compared
with treating $h$ and $t$ as standalone phrase nodes, this
representation keeps the relational triple intact, so a fact's role
in retrieval is determined by both its argument structure and the
passage it came from rather than by its endpoints alone.

\paragraph{Edges from shared endpoints and shared sources.}
The substrate links facts through two complementary signals derived
from the same indexing pass. \emph{Endpoint edges} connect two facts
$f, f'$ that share at least one informative argument phrase under
$\{h,t\}\cap\{h',t'\}\neq\emptyset$ when the phrases come from
different source passages, encoding that the same entity is described
across documents. \emph{Source edges} connect facts that share a
provenance, i.e.\ $\mathrm{src}(f)=\mathrm{src}(f')$, regardless of
endpoint overlap, encoding that they are co-asserted by the same
passage. Each fact node also carries a passage link to its provenance
$\mathrm{src}(f)$, so passage scores can later be read out by
aggregating mass on fact nodes grounded in $d$. We pre-compute dense
embeddings for fact units and relation phrases and materialize the
substrate as an adjacency index built once for $\mathcal{D}$; query
time performs no further OpenIE extraction over the corpus. Because
endpoint edges connect facts across passages while source edges
connect facts within a passage, the same substrate naturally supports
both cross-document propagation and local fact aggregation, which is
the property exploited by the query-local construction of
\S\ref{sec:method:retrieval}.

\subsection{Evidence-Driven Local Graph Retrieval}
\label{sec:method:retrieval}

\paragraph{Query-local subgraph.}
Given $q$, \methodname{} first induces a query-local subgraph $G_q$ that
selects the part of the offline fact--source substrate considered for
the query. We extract entity anchors
$A(q) \subseteq \mathcal{V}_{\mathrm{phr}}$ from $q$ using a controlled
language model that operates on the question only, never on retrieved
documents. A standard dense retriever provides an entry signal
$\mathrm{Entry}(q) \subseteq \mathcal{D}$, used as an activation prior
that seeds $G_q$ but does not by itself determine the final ranking.
The seed passage set is
\begin{equation}
\label{eq:gq-seed}
\begin{aligned}
\mathcal{S}_q
= {}& \mathrm{Entry}(q) \\
& {}\cup
\{\,\mathrm{src}(f) : f\in\mathcal{F},\\
&\qquad \{h,t\}\cap A(q)\neq\emptyset\,\}.
\end{aligned}
\end{equation}
From $\mathcal{S}_q$, $G_q$ is the bounded-hop expansion that follows
endpoint edges to facts in other passages and source edges within
each passage, for a fixed budget of hops over the substrate of
\S\ref{sec:method:substrate}. Thus $G_q$ contains the fact and passage
nodes that can participate in query-time retrieval, while excluding
the rest of the corpus graph, which keeps the activated region compact
yet still admits passages reachable through multi-hop entity chains
even when their surface similarity to $q$ is low.

\paragraph{Local personalized PageRank.}
The subgraph $G_q$ selects the activated region but does not itself
produce a passage ranking, since its fact and passage nodes play
different roles in retrieval. We therefore build a derived graph
$H_q$ on top of $G_q$ and run a personalized PageRank tailored to
the fact-source substrate, which differs from the corpus-level PPR
over passage-phrase graphs used in prior OpenIE-based retrieval.
The state space of $H_q$ is the query-local fact-and-passage layer
of $G_q$ rather than the entire corpus graph. Propagation moves
through OpenIE \emph{facts} rather than phrase nodes, traversing
multi-hop chains at the relational unit level. And the endpoint and
source edges of \S\ref{sec:method:substrate} act as two complementary
propagation modes within the same iteration. Endpoint edges carry
mass across passages, while source edges carry mass within a passage.
The resulting query-conditioned mass distribution $\pi_q$, which we
call the \emph{evidence flow}, is therefore both local to $q$ and
grounded in the offline graph.

Concretely, $H_q$ has two node types, local fact nodes and local
passage nodes. A fact $f=(h,r,t) \in \mathcal{F}$ belongs to
$\mathcal{F}_q$ when its provenance $\mathrm{src}(f)$ is a passage
node in $G_q$. Because the passage layer of $G_q$ is already
question-conditioned, $\mathcal{F}_q$ inherits this locality through
provenance, with no extra phrase-side condition. Two fact nodes are
linked by the endpoint and source edges of
\S\ref{sec:method:substrate} restricted to $\mathcal{F}_q$, and each
fact node links to its provenance passage. The seed distribution
$s_q$ places mass on facts whose arguments include an anchor in
$A(q)$, weighted by entry scores of the corresponding passages. We
solve, by power iteration on $H_q$,
\begin{equation}
\label{eq:local-ppr}
\pi_q = \alpha\, s_q + (1-\alpha)\, P_q^{\!\top}\, \pi_q,
\end{equation}
where $P_q$ is the transition matrix induced by fact--fact and
fact--passage links in $H_q$. The passage score for $d$ aggregates
stationary mass on local facts grounded in $d$:
\begin{equation}
\label{eq:passage-score}
s_{\mathrm{ppr}}(d \mid q)
=
|\mathcal{F}_{q,d}|^{-1/2}
\sum_{f \in \mathcal{F}_{q,d}} \pi_q(f).
\end{equation}
Here,
$\mathcal{F}_{q,d}=\{\,f\in\mathcal{F}_q:\mathrm{src}(f)=d\,\}$.
The square-root normalization keeps the passage score robust to
deduplication of the underlying fact set. Because propagation is
restricted to $H_q$, $s_{\mathrm{ppr}}$ rewards passages connected to
$q$ through entity-relation evidence rather than passages independently
similar to $q$.

\subsection{Evidence-Aware Retrieval Optimization}
\label{sec:method:enhancement}

The local PPR ranking $s_{\mathrm{ppr}}$ is graph-aware but operates
under the boundary of $G_q$ and treats each passage in isolation, so
two reader-facing problems remain. First, a passage on a multi-hop
chain can fall outside $G_q$ or sit several hops from any anchor,
receiving little $\pi_q$-mass even when its relation phrase exactly
matches what $q$ asks for. Second, the head of the ranking can crowd
around a single entity neighborhood, leaving other entities or
relations demanded by $q$ uncovered in any short prefix. We therefore
extend $s_{\mathrm{ppr}}$ with two evidence-aware steps that target
these problems directly. Relation-grounded expansion enlarges the
candidate pool with passages that the substrate links to $\pi_q$-mass
through a query-relevant relation, and coverage-aware reranking
selects a final ranking that addresses distinct question-side
requirements rather than repeating the strongest neighborhood.

\paragraph{Relation-grounded expansion.}
This step targets the bridge-underweighting problem by promoting
passages whose relation phrase carries query-relevant evidence,
including passages outside $G_q$. For a candidate fact
$e=(h_e, r_e, t_e)$ in $\mathcal{F}$ with provenance
$\mathrm{src}(e)=d$, we score $d$ by combining the largest
$\pi_q$-mass on any local fact $f=(h_f, r_f, t_f) \in \mathcal{F}_q$
sharing an argument phrase with $e$ and the dense compatibility
between $r_e$ and $q$,
\begin{equation}
\label{eq:expansion-score}
\begin{aligned}
s_{\mathrm{exp}}(d \mid q)
= {}& \max_{\substack{e:\,\mathrm{src}(e)=d \\ f \in \mathcal{F}_q}}
\Bigl\{\\[-0.3em]
&\pi_q(f)\,\bigl[\sigma(\mathbf{r}_e,\mathbf{q})\bigr]_{+}\\
&{}\cdot\mathbf{1}\!\left[
\{h_f,t_f\}\cap\{h_e,t_e\}\neq\emptyset
\right]\Bigr\},
\end{aligned}
\end{equation}
where $\mathbf{r}_e$ and $\mathbf{q}$ are dense embeddings of the
relation phrase and the question, $\sigma$ is cosine similarity, and
$[x]_{+}=\max(x,0)$ clips negative similarities so that expansion can
only add evidence. The base score then linearly combines local PPR
and expansion evidence,
\begin{equation}
\label{eq:base-score}
s_{\mathrm{base}}(d \mid q)
=
s_{\mathrm{ppr}}(d \mid q)
+
\eta\,s_{\mathrm{exp}}(d \mid q),
\end{equation}
with $\eta$ the only added scalar, tuned once on a development split.
Passages without an expansion path have
$s_{\mathrm{exp}}(d \mid q)=0$, so the base score reduces to local PPR
in the absence of relational evidence. Let $C_q$ denote the union of
passages scored by local PPR and passages admitted by expansion.

\paragraph{Coverage-aware reranking.}
This step targets the entity-neighborhood crowding problem by
selecting passages that collectively address the distinct
question-side requirements $B(q)=\{b_1,b_2,\ldots\}$ extracted from
$q$ by the same controlled language model used for $A(q)$. For a
passage $d$ and requirement $b$, a contribution
$\phi_{d,b}\in[0,1]$ is derived from dense similarity between $d$
and $b$ under $q$. Defining the noisy-OR coverage of a selected set
$R$ as $\mathrm{Cov}_q(R)=\sum_{b\in B(q)}\bigl[1-\prod_{d'\in R}(1-\phi_{d',b})\bigr]$,
which saturates once any selected passage substantially contributes
to $b$, the marginal coverage gain of adding $d$ is
\begin{equation}
\label{eq:marginal-gain}
\Delta_q(d \mid R) = \mathrm{Cov}_q(R\cup\{d\}) - \mathrm{Cov}_q(R).
\end{equation}
We then form $R_q$ by the greedy marginal-gain procedure of MMR-style set selection
\citep{carbonell1998mmr,lin2011submodular}, at each step adding the
candidate that maximizes
\begin{equation}
\label{eq:greedy-rerank}
d^{*}
=
\operatorname*{arg\,max}_{d \in C_q \setminus R}
\lambda\,s_{\mathrm{base}}(d \mid q)
+
(1-\lambda)\,\Delta_q(d \mid R),
\end{equation}
with $\lambda \in [0,1]$ the only added scalar, also tuned once on a
development split. Because per-requirement coverage saturates,
passages duplicating already-covered neighborhoods receive
diminishing marginal gain, while passages addressing previously
uncovered requirements are promoted. The final output is the ranked
evidence list $R_q$ over $\mathcal{D}$, and downstream systems may
consume any prefix of $R_q$.
